\section{Status of the Local Heat-Bath Approach}

Theorem~\ref{thm:dense-rank-recursion} gives a polynomial-time sampler for the
dense schedule without proving that a local chain mixes rapidly.  This section
records the still-distinct status of the exact local
$\{3,4,5\}$ heat-bath approach.  Its stationary law and several structural
inputs are known. Theorem~\ref{thm:balanced-bottleneck} now refutes a
universal polynomial mixing bound; rapid mixing in specified long-length
regimes remains open.

\subsection{Proved inputs}

The exact block heat-bath kernel is reversible for the uniform law on
$\Omega(N,m)$.  Therefore, for any fixed finite $(N,m)$, a lazy irreducible
version of the chain converges to the desired global uniform distribution.

Pure Hurwitz walks are confined to Hurwitz components and therefore cannot be
global samplers in the small cases of Table~\ref{tab:hurwitz-obstruction}.

The direct $r=4$ and $r=5$ conditional samplers are exact in the
all-transpositions model.  Exact block counts and per-block probabilities were
validated through $S_5$ for both $r=4$ and $r=5$, and the integrated
$\{3,4,5\}$ heat-bath transition checks passed $S_3,m=4$, $S_4,m=4$, and
$S_4,m=6$.

\begin{lemma}[Feasible identity components]
Let $C$ be a labeled vertex set with $|C|=v\ge 2$.  There exists an identity
word of length $k$ whose factors are transpositions supported on $C$ and whose
support graph is connected if and only if
\[
  k \text{ is even}\qquad\text{and}\qquad k\ge 2(v-1).
\]
\end{lemma}

\begin{proof}[Proof sketch]
The parity condition is forced because an identity permutation is even.  The
lower bound $k\ge 2(v-1)$ is the standard transitive transposition
factorization lower bound specialized to the identity: a transitive
factorization of a permutation $\sigma\in S_v$ has length at least
$v+c(\sigma)-2$, and $c(\id)=v$.  Conversely, choose any tree on $C$ and list
each of its $v-1$ edges as an adjacent identical pair.  The resulting word has product $\id$ and
connected support.  Additional length can be added in increments of two by
inserting square pairs $(t,t)$.
\end{proof}

The latest proof runs also isolated a useful macro-fiber structure.  Fix the
labeled weighted support-component data of a word: the nonzero support
components, together with the number of factors assigned to each component.
Since transpositions supported on disjoint components commute, the conditional
state space factors into a color-word interleaving problem and independent
transitive identity-factorization problems inside the components.  The
restricted heat-bath chain on such a fiber is reversible for the conditional
uniform law.  This factorization is now a proved structural input, but its
spectral consequences remain open.

\subsection{Unconditional irreducibility}

\begin{proposition}[Irreducibility]
The lazy exact $\{3,4,5\}$ same-product heat-bath chain is irreducible
on $\Omega(N,m)$ for $N\ge2$ and even $m\ge4$.
\end{proposition}
\begin{proof}
Lemma~\ref{lem:all-endpoint-pair-connectivity} connects every endpoint
fiber using adjacent pair replacements.  Every adjacent pair lies in
some nonwrapping length-three window when $m\ge3$.  The corresponding
pair replacement, with the third gate retained, has positive probability
under that window's exact heat bath.  Thus the project chain contains all
edges needed for connectivity.
\end{proof}

Earlier manuscript versions made this statement conditional on connected
identity Hurwitz transitivity.  The constructive pair-connectivity proof
removes that dependency.  No quantitative gap estimate follows from this
connectivity argument alone.

The exceptional cases are understood separately.  If $m$ is odd then
$\Omega(N,m)$ is empty.  If $m=0$ it is a singleton.  If $N=2$ and $m$ is even
it is also a singleton.  If $m=2$ and $N\ge 3$, then
$\Omega(N,2)=\{(t,t):t\in T_N\}$ is disconnected for a chain whose allowed
block sizes are $3,4,5$.

\subsection{Conditional polynomial-mixing statements}

A polynomial mixing bound over all $N,m$ is false by
Theorem~\ref{thm:balanced-bottleneck}. The following conditional routes
remain useful on restricted parameter regimes; their polynomial premises
cannot all hold uniformly across the bottleneck family.

For any lazy reversible irreducible chain $K$ on $\Omega(N,m)$ with spectral
gap $\gamma$,
\[
  \tau_{\mathrm{mix}}(\epsilon)
  \le
  \gamma^{-1}\left(m\log\binom{N}{2}+\log\frac{1}{\epsilon}\right),
\]
because the uniform stationary law has
$\pi_{\min}\ge |T_N|^{-m}=\binom{N}{2}^{-m}$.  Thus an inverse-polynomial lower bound on the gap of $\Klocal$ would
suffice on any specified regime where such a bound is possible.

\begin{theorem}[Projection/decomposition route, conditional]
Let $\Phi$ be the labeled weighted
support-component map, and let $\bar K$ be the stationary-flow projection of
$\Klocal$ onto the image of $\Phi$.  Suppose the following estimates hold:
\begin{enumerate}
\item a projection Poincare inequality with constant $q_{\mathrm{proj}}(N,m)$;
\item a restricted-fiber Poincare inequality with constant
  $q_{\mathrm{fib}}(N,m)$ for every fiber of $\Phi$;
\item a mean-lift bound with constant $q_{\mathrm{lift}}(N,m)$ controlling the
  projected Dirichlet form of fiber averages by the original Dirichlet form.
\end{enumerate}
Then
\[
  \mathrm{gap}(\Klocal)
  \ge
  \frac{1}{q_{\mathrm{fib}}(N,m)
    +q_{\mathrm{proj}}(N,m)q_{\mathrm{lift}}(N,m)}.
\]
In particular, if the three constants are polynomial in $N$ and $m$, then
$\Klocal$ mixes in time polynomial in $N$, $m$, and $\log(1/\epsilon)$.
\end{theorem}

\begin{proof}[Proof sketch]
Use the variance decomposition
$\mathrm{Var}_{\pi}(f)=\mathrm{Var}_{\bar\pi}(\mathbb{E}[f\mid\Phi])+
\mathbb{E}_{\bar\pi}[\mathrm{Var}(f\mid\Phi)]$.  Apply the projection
Poincare inequality to the first term, the restricted-fiber inequality to the
second term, and the mean-lift bound to compare the projected Dirichlet form
back to the Dirichlet form of $\Klocal$.
\end{proof}

\begin{theorem}[Large-block comparison route, conditional]
Let $K_L$ be the lazy ideal length-$L$ same-product
heat-bath chain.  Suppose that, for some polynomially bounded $L=L(N,m)$,
$K_L$ has inverse-polynomial gap, every positive-probability $K_L$ transition
can be simulated by $\{3,4,5\}$ subwindow updates, and the randomized
simulations have polynomial Diaconis-Saloff-Coste comparison congestion.  Then
$\Klocal$ has inverse-polynomial gap and hence polynomial mixing time.
\end{theorem}

\begin{proof}[Proof sketch]
Route the stationary flow of the reversible ideal chain $K_L$ through local
$\{3,4,5\}$ paths and apply the standard reversible-chain comparison theorem.
The gap loss is the product of the inverse gap of $K_L$ and the comparison
congestion, up to absolute laziness and normalization constants.
\end{proof}

\subsection{Current proof obligations}

The bottleneck cut is itself a union of support profiles, so the projection
cannot have an inverse-polynomial gap throughout the original parameter
range. On restricted regimes this route still needs
a projection conductance or Poincare bound;
the weighted-component fibers must mix in polynomial time; and the mean-lift
term must not lose more than a polynomial factor.  The latest run proved a
fresh-square certificate for support expansion from a length-$4$ identity
window and found no small support-size or weighted-component bottleneck through
$S_5,m=6$.  For example, the $S_5,m=6$ support-size conductance was about
$0.614126$, and singleton/threshold scans of the labeled weighted-component
projection had minimum conductance about $0.556800$.  These finite checks did not detect the asymptotic bottleneck. At long
lengths Theorem~\ref{thm:long-profile-conductance} supplies a profile
conductance bound for the pair chain, but leaves its restricted-fiber
conductance unestimated.

For the large-block route, $K_m$ samples the uniform bridge in one non-lazy
step, but this is too global to implement locally.  A useful mesoscopic
$K_L$, with $L<m$, would require parity-aware bridge smoothing and a
one-dimensional block-dynamics contraction theorem.  Local simulation is false
for $L=2$ in general.  For $L\ge 6$, the remaining requirement is a
quantitative pair-resampling and congestion problem.
Lemma~\ref{lem:all-endpoint-pair-connectivity} settles connectivity;
a useful gap comparison still needs bounds on path lengths and congestion.
Theorem~\ref{thm:balanced-bottleneck} prevents polynomial comparison
with a rapidly mixing ideal chain throughout the full parameter range.

The main negative theory lesson is that deterministic routing through one
canonical doubled-tree or square-rich normal form should not be used as a
rapid-mixing proof.  If positive comparison flow is forced through a fixed
canonical word $c$, vertex congestion is at least on the order of $1/\pi(c)$;
for a fixed canonical square word, $\pi(c)=1/|\Omega(N,m)|$.  Normal forms are
therefore useful for connectivity, but a mixing proof needs randomized routes
that stay in typical macrostates.
